\section{Why Consciousness Research Belongs Here}

The original framework treated behavioural states as practical objects only.
That was enough for an early conceptual draft, but it leaves the book too
detached from current biology. If conscious access, wakefulness, attention,
and large-scale coordination are themselves state-like biological phenomena,
then the language of state persistence and state transition should not stop at
everyday productivity.

This chapter therefore follows the manuscript's evidence policy explicitly. The
\textbf{strong empirical core} is the indexed set of official sources assigned
to this chapter in the literature map: \emph{Dynamical structure-function
correlations provide robust and generalizable signatures of consciousness in
humans} (Communications Biology, 2024),
\emph{Cortical connectivity, local dynamics and stability correlates of global
conscious states} (Communications Biology, 2025), and
\emph{Falling asleep follows a predictable bifurcation dynamic}
(Nature Neuroscience, 2025). The \textbf{medium} layer is interpretive:
classic guides such as Stanislas Dehaene's \emph{Consciousness and the Brain}
and Christof Koch's \emph{Consciousness} help explain why those findings
matter for this book, but they are not used as substitutes for the empirical
core. \textbf{Speculative} consciousness-physics proposals remain outside this
chapter and belong in Appendix~\ref{app:frontier}.
